\section{Estrategia de idempotencia}
\label{sec:idempotencia}

\subsection{El problema que ACID no resuelve}

El servidor confirma la transferencia y la respuesta se pierde antes de llegar
al cliente. El cliente no sabe si pagó. Si reintenta y el servidor trata la
petición como nueva, se paga dos veces. El problema ocurre \emph{fuera} del
límite de la transacción, así que ninguna propiedad ACID lo cubre.

\subsection{Diseño adoptado}

Tabla \codigo{idempotency\_keys} en PostgreSQL con unicidad sobre
\codigo{(user\_id, key)}, huella SHA-256 del método, la ruta y el cuerpo
canonicalizado, y un estado que distingue las peticiones en curso de las
finalizadas.

El protocolo tiene cuatro caminos, resumidos en el
cuadro~\ref{tab:idempotencia}.

\begin{table}[htbp]
\centering
\caption{Comportamiento ante una clave de idempotencia.}
\label{tab:idempotencia}
\small
\begin{tabularx}{\textwidth}{@{}lXl@{}}
\toprule
\textbf{Situación} & \textbf{Comportamiento} & \textbf{Respuesta} \\
\midrule
Clave nueva & Se reclama con \codigo{INSERT ... ON CONFLICT DO NOTHING} y se
  procesa & 201 o 422 \\
Misma clave, mismo cuerpo, finalizada & Se reproduce la respuesta guardada &
  la original, con \codigo{Idempotent-Replay: true} \\
Misma clave, mismo cuerpo, en curso & No se procesa nada & 409 \\
Misma clave, cuerpo distinto & No se procesa nada & 422 \\
\bottomrule
\end{tabularx}
\end{table}

Los códigos 409 y 422 son los que recomienda el borrador del IETF sobre la
cabecera \codigo{Idempotency-Key}~\cite{ietfidem}. Conviene precisar que ese
documento está \textbf{expirado}, por lo que se cita como práctica emergente y
no como estándar. El tiempo de vida de 24 horas está alineado con la política
publicada de Stripe para su API v1~\cite{stripeidem}.

\subsection{La decisión que hace que esto funcione}

La clave se reclama y se cierra \textbf{dentro de la misma transacción} que
mueve el dinero. Si se escribiera después del \codigo{COMMIT}, una caída entre
ambos pasos dejaría una transferencia ejecutada sin registro de idempotencia,
y el reintento la duplicaría. Al compartir transacción, la atomicidad del
motor cubre las dos cosas: o se movió el dinero y la clave quedó escrita, o no
ocurrió ninguna de las dos.

Esta es también la razón por la que Redis quedó descartado para este uso:
poner la clave y el dinero en sistemas distintos reabre exactamente la ventana
que el diseño cierra.

\subsection{Ámbito por usuario}

La unicidad es \codigo{(user\_id, key)} y no solo \codigo{key}. Evita
colisiones entre clientes y, más importante, impide que alguien reutilice la
clave de otro para obtener su respuesta almacenada. Es una propiedad de
seguridad además de una de corrección, y hay una prueba que la verifica.
\estado{verificado}

\subsection{Un caso que el diseño no había previsto}

Al implementar apareció una situación que el diseño sobre papel no
contemplaba. Si el \codigo{INSERT} choca con una fila que otra transacción aún
no ha confirmado, en Read Committed la consulta posterior \emph{tampoco la
ve}: no se reclama la clave y no se encuentra la existente.

Semánticamente es justo el caso de petición concurrente en curso, así que se
responde 409. Hay una prueba con dos hilos y dos conexiones reales que
comprueba que exactamente uno reclama la clave y que el dinero se mueve una
sola vez. \estado{verificado}

\subsection{El lado del cliente}

La interfaz web genera la clave una vez por intento de operación y la
\textbf{conserva entre reintentos}. Si la petición falla por red, no la
descarta: muestra \enquote{Resultado no confirmado}, cambia el botón a
\enquote{Reintentar envío} y reenvía la misma clave. De poco sirve un servidor
idempotente si el cliente genera una clave nueva en cada intento.
